\section{Waking, Dreaming, and the Wild Face}
\label{sec:dreams}

Much of the inner life lies on the wild side of the field, the part that has not settled into a
fixed, shared pattern, though the inner life has its lawful, settled structure too. Dreaming is
one window onto the wild side, and one thing the model can do is picture waking and dreaming as a
single mesh in two regimes, differing in whether the shared world is holding it.

\paragraph{One mesh, two regimes.}
Imagine a mesh that has settled several stable patterns into itself, the way a memory holds
several things it can recognize. In \emph{waking}, external input clamps part of the mesh to
what the senses present, and the mesh completes to the matching stored pattern and holds it,
steady and veridical, pinned to the shared world. It stays on the one pattern the world is
showing it. In \emph{dreaming}, that clamp is released. A slow internal rhythm then sweeps the
mesh through its own stored patterns, wandering its whole landscape and recombining what it
finds, with no external reality to anchor it. In the testbed this difference is just the
presence or absence of the external clamp: the same mesh, the same rule, the same memories,
pinned to reality when awake and free to roam its own depths when asleep.

\paragraph{What dreams might be.}
We should be careful here, because dreaming has not been examined deeply enough in this program
to say with any confidence what it is. One natural reading of the two regimes is that a dream is
the mesh roaming its own landscape with the world let go, so that it would be built largely from
one's own memories, recombined into arrangements that never occurred, with waking imagination a
milder, partly clamped version of the same thing. This reading fits the testbed picture and may
well capture part of what happens. But it is only one possibility, and the same substrate
suggests others. Because a self can in principle receive influence through the field
(Section~\ref{sec:spiritual}), a dream could be, in whole or in part, a stream of experience
reaching the self from elsewhere, from another self, another region of the field, or what the
later section treats as higher realms. Still other accounts remain open. So the released-mesh
picture is offered as a plausible account of some of what dreaming may involve, not as a claim
about what every dream is.

\paragraph{The wild side in general.}
More broadly, the wild is wherever the field stays fluid and unsettled rather than holding a
fixed pattern, in any domain. The physical world is a lawful region, the part that has frozen
into a pattern everyone shares, but lawful and wild are conditions, not the same as physical and
inner. Dreaming and imagining are mostly on the wild side, fluid and unmade, while still being
real structure in the same mesh, just not locked to the common view. Lawful and wild are one
ocean, part frozen into reliable ice and part still moving water, and sleep is in part a nightly
loosening of the ice.

\paragraph{The honest edge.}
The structural account is reasonably clear, and the difference between the clamped and the free
regime is something a model can show. Whether the dream is \emph{felt}, what it is actually like
from the inside, is the same hard-problem boundary that runs under the whole picture. The model
can describe the dreaming regime. It cannot, from the outside, hand over the dream.
